\section{A degree-two matching identity}\label{sec:matching}

Write
\[
 Q(y,z)=\tfrac12\bigl(Q(y+z)-Q(y)-Q(z)\bigr)
\]
for the symmetric bilinear polarization of $Q$, so that $Q(y,y)=Q(y)$.

For $S\in\mathbb S^3$ and $F,H\in\mathbb C^{3\times2}$, put
\[
 \langle F,H\rangle_S
 =\operatorname{Re}\operatorname{tr}(F^*SH)
 =\operatorname{tr}(\widetilde F^T S\widetilde H).
\]
This is a symmetric real bilinear form.  A complex trigonometric
polynomial of degree at most $j$ is a linear combination of
$e^{i\ell\theta}$ for $-j\leq\ell\leq j$.

The finite lift will test $L_\theta$ against trigonometric polynomial
vectors built from the columns of a matrix $\widetilde C(\theta)$.
A constant Gram matrix for $\widetilde C$ makes the coefficient of
$\mathcal X$ in the quadratic contribution of the top blocks independent
of $\theta$. Requiring
$\langle F_\theta,C\rangle_S=\operatorname{tr}(SQ(Y_\theta,z))$
makes the mixed contribution linear in $Y_\theta$.
The next lemma achieves both properties with entries of
degree at most two. The matrix $S$ in the pairing may be indefinite.

\begin{lemma}[Degree-two matching]\label{lem:matching}
Let $S\in\mathbb S^3$ and let $z\in\mathbb R^3$ satisfy
$(z_1,z_2)\ne(0,0)$.  There is a matrix
$C(\theta)\in\mathbb C^{3\times2}$ whose entries are trigonometric
polynomials of degree at most two such that, for all
$\theta,\eta,\xi\in\mathbb R$,
\begin{equation}\label{eq:matching}
 \begin{aligned}
 \widetilde C(\theta)\widetilde C(\theta)^T&=Q(z),\\
 \langle F_\theta(\eta,\xi),C(\theta)\rangle_S
 &=\operatorname{tr}\bigl(SQ(Y_\theta(\eta,\xi),z)\bigr).
 \end{aligned}
\end{equation}
\end{lemma}

\begin{proof}
Write $z=\tau d(\alpha)+\rho k$, where $\tau>0$, and set
\[
 w_0=e^{i\alpha},\qquad R_0=\tau U(\alpha)+\rho V_0,
 \qquad N=SR_0.
\]
By Lemma~\ref{lem:gram}, $\widetilde R_0\widetilde R_0^T=Q(z)$.
We will set $C(\theta)=R_0M(\theta)$ with $M(\theta)$ unitary, so
that this Gram matrix remains fixed.  We choose the rotation to make
$\langle V_0,C\rangle_S$ constant and
$\langle U(\theta),C(\theta)\rangle_S$ a linear combination of
$\cos\theta$ and $\sin\theta$.  After those frequency bounds are
established, values and a derivative at $\alpha$ will identify the
pairings.

\smallskip
\noindent\emph{Choosing the unitary rotation.}
Write $N_j$ for the $j$th row of $N$ and introduce the two-component row
\[
 v=N_3-iN_2.
\]
Subscripts $+$ and $-$ denote the first and second components of a row.
Set
\[
 g=v_++\overline{v_-},\qquad
 h=v_+-\overline{v_-},\qquad
 \kappa(\theta)=(w_0/w)^2,\qquad
 \delta=|g|^2+|h|^2.
\]
We seek a matrix of the form
\[
 M(\theta)=\begin{pmatrix}a&b\\-\overline b&\overline a\end{pmatrix}.
\]
Writing $v^M=vM$, we impose the two conditions
\begin{equation}\label{eq:rotation-pairings}
 v^M_++\overline{v^M_-}=g,
 \qquad
 v^M_+-\overline{v^M_-}=h\kappa.
\end{equation}
The sum on the left controls the pairing with $V_0$, while the difference
multiplies $w^3$ in the pairing with $U$.  These conditions therefore
keep the first pairing constant and lower the cubic term of the second
to frequency one, since $w^3\kappa=w_0^2w$.  The expansion below will
also bound its remaining frequencies.

The two left sides of \eqref{eq:rotation-pairings} are
$ga+\overline h\,\overline b$ and
$ha-\overline g\,\overline b$, respectively.  Thus, for $\delta>0$,
we choose $a(\theta)$ and $b(\theta)$ by solving
\begin{equation}\label{eq:rotation-system}
 \begin{pmatrix}g&\overline h\\h&-\overline g\end{pmatrix}
 \begin{pmatrix}a\\\overline b\end{pmatrix}
 =\begin{pmatrix}g\\h\kappa\end{pmatrix}.
\end{equation}
The coefficient matrix $A$ satisfies $A^*A=\delta I_2$, so the
system has a unique solution.  Explicitly,
\begin{equation}\label{eq:rotation-coefficients}
 a=\frac{|g|^2+|h|^2\kappa}{\delta},\qquad
 \overline b=\frac{gh(1-\kappa)}{\delta}.
\end{equation}
Since the right side of \eqref{eq:rotation-system} has squared norm
$\delta$, we have $|a|^2+|b|^2=1$.
For $\delta=0$, take $a=1$ and $b=0$; then $g=h=0$ and the same
system is satisfied.  In both cases $M$ is unitary,
$M(\alpha)=I_2$, and $a,\overline b$ involve only the
frequencies $0$ and $-2$.  Thus $C=R_0M$ has degree at most two and
$C(\alpha)=R_0$.  Moreover,
\[
 \widetilde C\widetilde C^T
 =\operatorname{Re}(R_0MM^*R_0^*)
 =\operatorname{Re}(R_0R_0^*)=Q(z),
\]
which proves the first identity in \eqref{eq:matching}.

\smallskip
\noindent\emph{Reducing the frequencies of the pairings.}
The formula for $V_0$ and \eqref{eq:rotation-pairings} yield
\[
 \sqrt2\,\langle V_0,C\rangle_S
 =\operatorname{Re}(v^M_++v^M_-)=\operatorname{Re}g,
\]
so this pairing is constant.

Set $m=N_3+iN_2$, $n=N_1$, and write $m^M=mM$, $n^M=nM$.
Expanding $c$ and $s$ in powers of $w$, the formula for $U$ gives
\begin{align*}
 \sqrt2\,\langle U(\theta),C(\theta)\rangle_S
 =\operatorname{Re}\bigl[{}
 &\tfrac12w^3v^M_+
 +w(\tfrac12m^M_+-n^M_+)\\
 &-\tfrac12w^{-3}v^M_-
 +w^{-1}(-\tfrac12m^M_--n^M_-)
 \bigr].
\end{align*}
Inside the real part, we may replace each negative-power term by its
complex conjugate.  The resulting expression is
\begin{equation}\label{eq:frequency-reduction}
 \operatorname{Re}\left[
 \frac{w^3}{2}(v^M_+-\overline{v^M_-})
 +w\left(\frac{m^M_+}{2}-n^M_+
          -\frac{\overline{m^M_-}}{2}-\overline{n^M_-}\right)
 \right].
\end{equation}
The first term has frequency one because
$w^3\kappa=w_0^2w$.  In the second term, the expression in parentheses
has only frequencies $0$ and $-2$: for any fixed row $r$,
\[
 (rM)_+=r_+a-r_-\overline b,
 \qquad
 \overline{(rM)_-}=\overline{r_+}\,\overline b+\overline{r_-}a.
\]
Consequently \eqref{eq:frequency-reduction} is a real linear
combination of $\cos\theta$ and $\sin\theta$, as required.

\smallskip
\noindent\emph{Identifying the pairings.}
At $\theta=\alpha$, polarization of Lemma~\ref{lem:gram} within the
plane spanned by $d(\alpha)$ and $k$ gives
\begin{align}
 \langle V_0,R_0\rangle_S
 &=\operatorname{tr}\bigl(SQ(k,z)\bigr),\label{eq:matching-value-v}\\
 \langle U(\alpha),R_0\rangle_S
 &=\operatorname{tr}\bigl(SQ(d(\alpha),z)\bigr).
 \label{eq:matching-value-u}
\end{align}
Both sides of the desired identity
$\langle V_0,C(\theta)\rangle_S=\operatorname{tr}(SQ(k,z))$
are constant, so \eqref{eq:matching-value-v} proves it for every
$\theta$.

For the pairing with $U$, both
$\langle U(\theta),C(\theta)\rangle_S$ and
$\operatorname{tr}(SQ(d(\theta),z))$ are linear combinations of
$\cos\theta$ and $\sin\theta$.  They agree at $\alpha$ by
\eqref{eq:matching-value-u}; it remains to compare their derivatives
there.  The fixed Gram matrix of $C$ gives
\[
 \langle C(\theta),C(\theta)\rangle_S
 =\operatorname{tr}(SQ(z)),
 \qquad
 \langle R_0,C'(\alpha)\rangle_S=0.
\]
Constancy of the pairing with $V_0$ also gives
$\langle V_0,C'(\alpha)\rangle_S=0$.  Since
$R_0=\tau U(\alpha)+\rho V_0$ and $\tau>0$, it follows that
\begin{equation}\label{eq:matching-moving-factor}
 \langle U(\alpha),C'(\alpha)\rangle_S=0.
\end{equation}
Next differentiate the identity
\[
 \langle\tau U(\theta)+\rho V_0,
          \tau U(\theta)+\rho V_0\rangle_S
 =\operatorname{tr}\bigl(SQ(\tau d(\theta)+\rho k)\bigr),
\]
which follows from Lemma~\ref{lem:gram}.  Evaluating at $\alpha$ and
dividing by $2\tau$ gives
\[
 \langle U'(\alpha),R_0\rangle_S
 =\operatorname{tr}\bigl(SQ(d'(\alpha),z)\bigr).
\]
Together with \eqref{eq:matching-moving-factor}, this is the required
derivative identity.  A real linear combination of $\cos\theta$ and
$\sin\theta$ is determined by its value and derivative at a single
angle.  Hence the pairing with $U$ has the desired value for all
$\theta$.  Linearity in $\eta,\xi$ completes the proof.
\end{proof}
